\begin{lemma}[Change lemma, mixed location: no closed starred stage]
Fix a finite loopless multigraph with its exact edge set, a palette of
$k$ colors, and a distinguished edge $e_0$.  Let $(T,\varphi_n)$ be a
positive-rung extended Tashkinov tree.  Let $\mathcal B$ be a hierarchy
base and let
\[
 T_n=T_{n,0}\subset T_{n,1}\subset\cdots\subset
 T_{n,q}\subset T_{n,q+1}=T
\]
be a good hierarchy under $\varphi_n$, where $q$ is its last divider
index and
\[
 A=T_{n,q}^{*}=\noderef{HierarchyStarredStage}(H,q).
\]
Assume that $T$ has maximum property, that $V(T)$ is not elementary,
that $\varphi_n$ is proper and leaves exactly $e_0$ uncolored, and that
the rung number is positive.  These data form a
\noderef{HierarchyCounterexample} $C$.

Assume the six lower-rung assertions of the technical theorem, uniformly
for every positive-rung maximum ETT on the same graph, palette, and
distinguished edge, with arbitrary certified input and generating
colorings and with rung number smaller than that of $T$: lower-rung
elementarity, terminal interchangeability, the two PE stability and
stored-pair assertions, the terminal PE component assertion, the
defective-endpoint order assertion, and the mod-coloring assertion.
Assume also the same-rung induction hypothesis that every maximum ETT
with a good hierarchy having fewer than $q$ dividers is elementary.
Finally assume the large-color estimate
\[
 2n+11\leq |\overline\varphi_n(T_n)|.
\]

Put $\ell=|T|-|A|$, the number of exact tail additions above $A$, and let
\[
 T(y_{\ell-1})=T\big|_{\min\{|A|+(\ell-1),|T|\}}
\]
be the exact ordered prefix supplied by \noderef{TreeSequencePrefix}.
Assume (6.6): $V(T(y_{\ell-1}))$ is elementary under $\varphi_n$.
Assume moreover that $C$ is lexicographically minimal among all
\noderef{HierarchyCounterexample}s on the fixed graph, palette, and
$e_0$ which have exactly $q$ dividers, allowing their certified input
colorings, generating colorings, bases, ETTs, and good hierarchies to
vary.  Explicitly, no such counterexample $C'$ satisfies
\[
 \noderef{HierarchyCounterexampleLess}
  \bigl(\noderef{HierarchyCounterexampleMeasure}(C'),
        \noderef{HierarchyCounterexampleMeasure}(C)\bigr).
\]

Let $\alpha\ne\beta$ be colors missing on $T(y_{\ell-1})$.  Write
\[
 i_c=\noderef{HierarchyFirstMissingPosition}(\varphi_n,T,c),
 \qquad v_c=T.\mathit{vertex}(i_c).
\]
If at least one of $\alpha,\beta$ is not missing on $A$, assume both
\[
 v_\alpha\prec v_\beta                                      \tag{a}
\]
and the exact edge-identity exclusion
\[
 \alpha\notin
 \left\{c:\exists t\in\operatorname{Fin}|T|,
       |A|\leq t<i_\beta,
       \ \varphi_n(T.\mathit{edge}(t))=c\right\}.             \tag{b}
\]
Thus (b) is precisely
$\alpha\notin\varphi_n\langle T(v_\beta)-A\rangle$; it is not
an endpoint-only abbreviation.


Assume in addition that the first $\alpha$-missing vertex lies in the
last starred stage and the first $\beta$-missing vertex lies outside it:
\[
        v_\alpha\in V(A),\qquad v_\beta\notin V(A).
\]
Assume that $q>0$ and that, for every $s$ with
$0\leq s\leq q$, the singleton $\{\alpha\}$ is not closed on
$T_{n,s}^{*}$.  Then
\[
 \noderef{TwoColorComponent}(G,\varphi_n,\alpha,\beta,
       \operatorname{inl}v_\alpha)
 =
 \noderef{TwoColorComponent}(G,\varphi_n,\alpha,\beta,
       \operatorname{inl}v_\beta).
\]
\end{lemma}
\begin{proof}
Put $H=C.\mathit{hierarchy}$, $A=T_{n,q}^{*}$, and for each vertex $x$
write
\[
 K_x=\noderef{TwoColorComponent}
       (G,\varphi_n,\alpha,\beta,\operatorname{inl}x).
\]
The nested and terminal fields of \noderef{TashkinovHierarchy} make
$A$ and all its hierarchy stages exact ordered initial segments of
$T$.  The full tail length $\ell=|T|-|A|$ is positive.  Indeed, at
$\ell=0$, \noderef{TreeSequenceInitialSegment} gives $A=T$ and the
prefix \noderef{TreeSequencePrefix} in (6.6) is the full $T$, contrary
to the non-elementarity field of \noderef{HierarchyCounterexample}.
Consequently all hierarchy stages lie in the elementary prefix
$T(y_{\ell-1})$ and are elementary by restriction of
\noderef{ColoredVertexSetElementary}.

By \noderef{HierarchyFirstMissingPosition}, $v_\alpha$ misses $\alpha$
and $v_\beta$ misses $\beta$.  Since $q>0$,
\noderef{HierarchyStarredStage} identifies $A=T_{n,q}$; hence
$\alpha$ is missing on $T_{n,q}$.  The mixed location makes the two
vertices distinct, and elementarity of the full-tail prefix says that
neither misses the other color.  Thus \noderef{MissingColors},
\noderef{PresentColors}, and properness make $K_{v_\alpha}$ and
$K_{v_\beta}$ nontrivial path components with their individual edge
nodes retained.

Apply the second alternative of \noderef{rutcor}.  Its hypotheses are
exactly the present full-tail (6.6), maximum property, hierarchy, and
pre-step elementarity; the latter comes from lower assertion (i) after
restriction through the preceding rung, or from
\noderef{TechnicalTheoremBase} at rung one.  The conclusion gives all
four facts
\[
\begin{split}
 &\alpha\in
   \bigcup_{\eta\in D_{n,j}}\Gamma^{j-1}_\eta
       \subseteq\Gamma^{j-1}\quad(1\leq j\leq q),              \tag{1}\\
 &\Theta_n=PE,\qquad
 v_\alpha\in V(T_n)\setminus V(R_n),                           \tag{2}\\
 &\alpha\notin
      \varphi_n\langle T_{n,q}-T_n\rangle.                    \tag{3}
\end{split}
\]
In particular, this proof uses the reserved-color conclusion at every
stage, not merely the fact that the stages are nonclosed.

Suppose for contradiction that $K_{v_\alpha}\ne K_{v_\beta}$, and set
$\sigma_n=\varphi_n/K_{v_\beta}$ by
\noderef{TwoColorComponentSwap}.  The swap is proper by
\noderef{TwoColorComponentSwapProper}, preserves the unique uncolored
edge $e_0$, leaves $v_\alpha$ missing $\alpha$, and makes $v_\beta$
missing $\alpha$.

By conclusion (ii) of \noderef{interchange}, $K_{v_\alpha}$ is the
unique $(\alpha,\beta)$ path component meeting $T_n$; hence the distinct
component $K_{v_\beta}$ is disjoint from $T_n$.  No vertex of $R_n$
misses $\alpha$: otherwise it and
$v_\alpha\in T_n\setminus R_n$ would be two vertices of the elementary
core missing the same color.  No vertex of $R_n$ misses $\beta$ because
$v_\beta$ is the first vertex missing $\beta$ and lies outside $A$.
Thus at every vertex where $K_{v_\beta}$ meets $R_n$, both colors are
present and the switch merely exchanges them.  It follows that the
switch is stable on $T_n$ with control $D_n$, and separately stable on
$R_n$ with empty control.  These are exactly the two clauses of
\[
 \noderef{CombinedStableColoring}(T_n,R_n,D_n,\varphi_n,\sigma_n).
                                                                  \tag{4}
\]
This derives combined stability pointwise; PE alone is not being used
as a closure assertion.

Let $B=T(v_\beta)$ be the exact ordered prefix ending at $v_\beta$.
Writing $s_B=|B|-|A|$, the mixed location and
$v_\beta\in T(y_{\ell-1})$ give
\[
                       1\leq s_B\leq\ell-1<\ell.              \tag{5}
\]
The two vertices of $B$ now missing $\alpha$ make $B$ non-elementary
under $\sigma_n$.

Construct the old-color ETT on $B$ first.  Restrict the final
\noderef{TashkinovClosure} of the original ETT at the displayed length
$|B|$ while retaining its ladder series, exact connecting edges,
extension kinds, and maximum-property data.  Restrict only the terminal
hierarchy stage to $B$, keeping the base, all divider positions,
ordinary and starred stages, and Gamma pairs.  Since $A\prec B\prec T$,
the required terminal containments are strict.  Denote this positive
maximum ETT and hierarchy by $\mathcal E_B$.

Combine (3) with the ordered-unused condition beyond $A$ to exclude
$\alpha$ from every exact displayed edge of $B-T_n$.  The exact
closure sequence from $T_n$ to $R_n$ cannot use $\alpha$: such an
augmentation would require another vertex of the elementary core
missing $\alpha$, whereas (2) names its unique such vertex.  The same
first-missing argument excludes $\beta$ from the closure edges and all
ordinary edges before $v_\beta$.  Apply
\noderef{CombinedStableClosurePreserved} to (4).  It fixes every
displayed identity of $T_n$ and $R_n$ and transports the same exact
closure.  Together with the exclusions just proved, every edge in the
PE seed and subsequent augmentation route to $B$ keeps its old color.
Every missing-color witness is also preserved pointwise: outside the
switched component it is fixed, and on the component before
$v_\beta$ both $\alpha$ and $\beta$ are present and merely exchanged.
The restricted route is certified by
\noderef{TashkinovAugmentationSequenceRestrict} and consists of the
same \noderef{TashkinovAugmentation}s.

Apply \noderef{Td} to the $T_n$ stability clause of (4), using the
given six lower-rung assertions, and obtain \noderef{ModColoring}.
Then apply \noderef{HierarchySplitter} to the old-color prefix ETT
$\mathcal E_B$, using (4), the transported closure, and the verified
exact augmentation route.  It supplies a
\noderef{CorrespondingExtendedTashkinovTree} on the actual prefix $B$,
maximum property, a corresponding base, and
\noderef{SameHierarchyStages}.  In particular, it does not supply a
good hierarchy.

We construct that \noderef{GoodHierarchy} now, field by field, with
the old Gamma pairs.  At each stage vertex in the switched component,
both colors are present: $v_\alpha$ is in the other component and
$v_\beta$ occurs after $A$.  Therefore every stage missing-color set
is unchanged pointwise.  Pair size and pairwise disjointness are
literal.  Avoidance holds because the displayed segment edges were
shown above to retain their exact colors.  Unique refresh follows from
the unchanged pointwise missing sets on the exact stage differences.

For the PE parallel-base field, (4) preserves the missing sets on both
constituents, while
\noderef{CombinedStableClosurePreserved} preserves the exact closure;
hence the common-missing cardinality clause is unchanged.  To preserve
non-minus-closedness, take its old witness color and exact boundary
edge.  A changed witness color would be $\alpha$ or $\beta$.  At
$j=1$, (1) puts $\alpha$ in the reserved union $\Gamma^0$, so it cannot
be a witness from the missing colors outside $\Gamma^0$; and
$\beta$ is not missing on the core because its first missing vertex
lies outside $A$.  Therefore the witness edge is fixed and remains a
witness under $\sigma_n$.

At every positive stage $j$, (1) puts $\alpha$ in precisely the
reserved union excluded from the minus-closed demand, while $\beta$
is not missing there.  Thus any boundary edge changed from one of
these two colors to the other is irrelevant to the demanded color set;
all other boundary colors are fixed.  The positive-stage minus-closed
field follows for every $j$.  This completes all five good-hierarchy
fields.  Only after this construction do we invoke
\noderef{SameGoodHierarchy} to record the pointwise equality of stages,
stored sets, and Gamma pairs.

The corresponding maximum ETT on $B$, its base, hierarchy, complete
good witness, and the two vertices missing $\alpha$ form a
\noderef{HierarchyCounterexample} $C'$ with the same divider count.
Its measure is strictly smaller.  Let
$p=\noderef{HierarchyPathNumber}(H)$ and let $p'$ denote the truncated
path number.  If $s_B<p$, then $p'\leq s_B<p$.  If $p\leq s_B$, the
old path tail beginning at $p$ restricts to a path tail of $B$, with
the singleton case included, so $p'\leq p$.  If $p'=p$, set
$d=|A|-|T_{n,q}|$; then
\[
 |T|-|T_{n,q}|=d+\ell,\qquad
 |B|-|T_{n,q}|=d+s_B<d+\ell
\]
by (5).  Hence \noderef{HierarchyCounterexampleLess} holds between the
two \noderef{HierarchyCounterexampleMeasure}s, contradicting the
minimality of $C$.  Therefore the two original components are equal.
\end{proof}

